\chapter{Backend implementation} \label{ch:backend}

\section{Simulator module}

\subsection{\texttt{ArduinoState} class}

\begin{lstlisting}[language=csharp, caption={\texttt{ArduinoState} class}]
class $ArduinoState$ {
   public:
    static constexpr uint8_t ANALOG_PIN_COUNT = 6;
    static constexpr uint8_t DIGITAL_PIN_COUNT = 14;
    static constexpr uint8_t PIN_COUNT = 
        ANALOG_PIN_COUNT + DIGITAL_PIN_COUNT;

    using $analog_arr_t$ = $std$::$array$<$analog_t$, ANALOG_PIN_COUNT>;
    using $digital_arr_t$ = $std$::$array$<$digital_t$, DIGITAL_PIN_COUNT>;
    using $pinmode_arr_t$ = $std$::$array$<$PinMode$, PIN_COUNT>;

    void #init_timer#();
    float #get_time#() { return _timer.#deltaTime#(); }

    $analog_arr_t$& #get_analog#() { return _analogPins; }
    double #get_analog#($pin_t$ pin) const;
    bool #set_analog#($pin_t$ pin, $analog_t$ value);

    $digital_arr_t$& #get_digital#() { return _digitalPins; }
    bool #get_digital#($pin_t$ pin) const;
    bool #set_digital#($pin_t$ pin, $digital_t$ value);

    $pinmode_arr_t$& #get_pin_mode#() { return _pinModes; }
    $PinMode$ #get_pin_mode#($pin_t$ pin) const;
    bool #set_pin_mode#($pin_t$ pin, $PinMode$ mode);
   private:
    $Timer$ _timer;
    $analog_arr_t$ _analogPins = {};
    $digital_arr_t$ _digitalPins = {};
    $pinmode_arr_t$ _pinModes = {};
};
\end{lstlisting}

\subsection{\texttt{Simulator} singleton class}

\begin{lstlisting}[language=csharp, caption={\texttt{Simulator} class}]
class $Simulator$ {
   public:
    using $event_callback_fnct$ = 
        $std$::$function$<void(const $Event$&)>;
    using $logger_ptr$ = $std$::$shared_ptr$<logger::$ILogger$>;

    static void #init#(
        $event_callback_fnct$ eventCallback, 
        $logger_ptr$ logger = nullptr);

    static $ArduinoState$& #state#() { return s_instance->_state; }
    static $EventQueue$& #queue#() { return s_instance->_eventQueue; }

    static int64_t #simulation_time#() { return #state#().#get_time#(); }
    static int64_t #real_time#() 
        { return s_instance->_timer.#deltaTime#(); }

    static void #handle_event#($Event$ event);
    static void #handle_events#();

    static void #log_debug#($std$::$string$&& str) 
        { s_instance->_logger->#log_debug#($std$::#move#(str)); }
    static void #log_info#($std$::$string$&& str) 
        { s_instance->_logger->#log_info#($std$::#move#(str)); }
    static void #log_warning#($std$::$string$&& str) 
        { s_instance->_logger->#log_warning#($std$::#move#(str)); }
    static void #log_error#($std$::$string$&& str) 
        { s_instance->_logger->#log_error#($std$::#move#(str)); }
   private:
    static $Simulator$* s_instance;
    $event_callback_fnct$ f_eventCallback;
    $logger_ptr$ _logger;
    $ArduinoState$ _state;
    $EventQueue$ _eventQueue;
    $Timer$ _timer;
};
\end{lstlisting}

\subsection{\texttt{Event} struct}

\begin{lstlisting}[language=csharp, caption={\texttt{Event} struct}]
struct $Event$ {
   public:
    enum class $Type$ { Write, PinMode, Reboot };

    $Type$ type;
    int64_t timestampMicros;
    int64_t localVirtualTime;
    $std$::$function$<void()> action;

    static constexpr int MAX_ARGS = 3;
    $std$::$array$<int, MAX_ARGS> args;

    template <typename... Args>
    $Event$($Type$ type, Args... arguments);
    $Event$() {}

    size_t #id#() const { return _id; }
    $std$::$string$ #to_string#();

    static $Event$ #write#($pin_t$ pin, $digital_t$ value);
    static $Event$ #set_pinmode#($pin_t$ pin, $PinMode$ mode);
    static $Event$ #reboot#();
   private:
    static size_t s_nextEventId;
    size_t _id;
};
\end{lstlisting}

\subsection{\texttt{EventQueue} class}

\begin{lstlisting}[language=csharp, caption={\texttt{EventQueue} class}]
class $EventQueue$ {
   public:
    void #enqueue_local#($Event$& event);
    void #enqueue_remote#($Event$ event);
    void #execute_next_event#();

    $Event$& #last_executed_event#() 
        { return _processedEvents.#back#(); }
    bool #is_empty#() const 
        { return _localEvents.#empty#() && _remoteEvents.#empty#(); }
    int64_t #next_lvt#() const { return _nextLvt; }
    int64_t #last_lvt#() const { return _lastLvt; }

   private:
    $std$::$queue$<$Event$> _localEvents;
    $std$::$queue$<$Event$> _remoteEvents;
    $std$::$queue$<$Event$> _processedEvents;
    int64_t _nextLvt;
    int64_t _lastLvt;
    int64_t _maxGvtFallback;

    void #clear#();
    size_t #earliest_queue_lvt#(const $std$::$queue$<$Event$>& queue);
};
\end{lstlisting}

\subsection{\texttt{Arduino.cpp}}

\section{Executable module}

\subsection{\texttt{main.cpp}}

\begin{lstlisting}[language=csharp, caption={\texttt{main.cpp}}]
using $logger_t$ = logger::$FileLogger$;
using namespace $arguino$::$simulator$;
$ProgramOptions$ options;

void #init_simulator#(
    $Simulator$::$event_callback_fnct$ eventCallback)
{
    auto logPath = $std$::$filesystem$::#absolute#(
        options.SimulatorLogPath);
    auto logger = $std$::#make_shared#<$logger_t$>(logPath);

    $Simulator$::#init#(eventCallback, logger);
    logger->timerFunction = $Simulator$::#simulation_time#;
    logger->verbosityLevel = options.Verbosity;
}

void #simulator_loop#()
{
    $Simulator$::#state#().#init_timer#();

    $Simulator$::#handle_event#($Event$::#reboot#());
    $Simulator$::#log_info#("Simulator initialized.");
    #setup#();
    $Simulator$::#log_info#(
        "setup() finished successfully. Running loop().");

    while (true) {
        #loop#();
    }
}

int #main#(int argc, char** argv)
{
    options = #parse_arguments#(argc, argv);
    #init_simulator#(on_event);
    #run_simulator_with_ipc#(#simulator_loop#, options);
}
\end{lstlisting}

\subsection{\texttt{ipcAdapter.hpp}}

\begin{lstlisting}[language=csharp, caption={\texttt{ipc\_adapter.hpp}}]
using $logger_t$ = logger::$FileLogger$;
using $loop_fnct$ = $std$::$function$<void()>;
using $init_fnct$ = $std$::$function$<void(
    arguino::simulator::$Simulator$::$event_callback_fnct$)>;

static $std$::$shared_ptr$<$logger_t$> IpcLogger;

void #run_simulator_with_ipc#(
    $loop_fnct$ loopFunction, 
    const $ProgramOptions$& options);

void #on_event#(const arguino::simulator::$Event$& event);

inline void #init_ipc_logger#(const $ProgramOptions$& options)
{
    auto loggerAbsolutePath = $std$::$filesystem$::#absolute#(
        options.IcpLogPath);
    IpcLogger = $std$::#make_shared#<$logger_t$>(loggerAbsolutePath);

    IpcLogger->verbosityLevel = options.Verbosity;
    IpcLogger->timerFunction = 
        arguino::simulator::$Simulator$::#real_time#;
}
\end{lstlisting}

\subsection{\texttt{tcpAdapter.cpp}}

\begin{lstlisting}[language=csharp, caption={\texttt{tcp\_adapter.cpp}}]
using namespace arguino::simulator;
using $tcp_server_t$ = arguino::tcp::$TcpServer$;
static $std$::$unique_ptr$<$tcp_server_t$> _tcpServer;

void #on_event#(const $Event$& event)
{
    auto message = #encode_event#(event);
    if (message == UNKNOWN_EVENT) {
        IpcLogger->#log_error#(
            "Encountered unknown event type; skipping.");
        return;
    }

    IpcLogger->#log_debug#("Enqueuing message " + message);
    _tcpServer->#post_message#(message);
}

static void #on_received_tcp_message#(const $std$::$string$ message)
{
    $Event$ event;
    if (#decode_event#(message, event)) {
        $Simulator$::#queue#().#enqueue_remote#(event);
    }
    else {
        IpcLogger->#log_error#(
            "Could not decode event from message " + message);
    }
}

static void #run_tcp#() { _tcpServer->#launch#(); }

void #run_simulator_with_ipc#(
    $loop_fnct$ loopFunction, 
    const $ProgramOptions$& options)
{
    #init_ipc_logger#(options);
    _tcpServer = $std$::#make_unique#<$tcp_server_t$>(
        options.TcpPort, on_received_tcp_message, IpcLogger);

    $std$::$thread$ tcp_thread(#run_tcp#);
    IpcLogger->#log_info#("Tcp Server launched");
    loopFunction();
    tcp_thread.#join#();
}
\end{lstlisting}

\subsection{\texttt{shmemAdapter.cpp}}

\begin{lstlisting}[language=csharp, caption={\texttt{shmem\_adapter.cpp}}]
using namespace arguino::shmem;
using namespace arguino::simulator;

constexpr uint8_t MESSAGE_DELIMETER = ';';
static $TwoWayBuffer$* _shmem;

void #on_event#(const $Event$& event)
{
    auto message = #encode_event#(event);
    if (message == UNKNOWN_EVENT) {
        IpcLogger->#log_error#(
            "Encountered unknown event type while encoding.");
        return;
    }

    if (!_shmem->#write#(message + ';')) {
        IpcLogger->#log_error#(
            "Could not write to shared memory buffer, "
            "not enough available space.");
    }
    else {
        IpcLogger->#log_debug#(
            "Written message to producer buffer: " + message);
    }
}

static void #run_shmem_consumer#()
{
    while (true) {
        while (!_shmem->#consumer#().#is_empty#()) {
            $std$::$vector$<uint8_t> data = 
                _shmem->#consumer#().#consume_until#(
                    MESSAGE_DELIMETER);
            if (data.#empty#()) {
                IpcLogger->#log_error#(
                    "A span of incoming message could not "
                    "be determined, skipping.");
                continue;
            }

            $std$::$string$ message(data.#begin#(), data.#end#());
            IpcLogger->#log_debug#(
                "Read message from consumer buffer: " + message);

            $Event$ event;
            if (#decode_event#(message, event)) {
                $Simulator$::#queue#().#enqueue_remote#(event);
            }
            else {
                IpcLogger->#log_error#(
                    "Could not decode event from message " + message);
            }
        }
    }
}

void #run_simulator_with_ipc#(
    $loop_fnct$ loopFunction, 
    const $ProgramOptions$& options)
{
    #init_ipc_logger#(options);
    _shmem = new $TwoWayBuffer$(
        options.ShmemName, options.ShmemSizePages);

    $std$::$thread$ consumer_thread(#run_shmem_consumer#);
    IpcLogger->#log_info#("Shared memory buffer intitialized.");

    loopFunction();
    consumer_thread.#join#();
}
\end{lstlisting}

\subsection{\texttt{textEncoder.cpp}}

\begin{lstlisting}[language=csharp, caption={\texttt{encoder.cpp}}]
using namespace arguino::simulator;

constexpr char PART_DELIMETER = ':';
constexpr char COMMON_PART_FORMAT[] = "{:012}:{:07}:";
constexpr char WRITE_FORMAT[] = "W:{:02}:{}";
constexpr char PINMODE_FORMAT[] = "P:{:02}:{}";
constexpr char REBOOT_FORMAT[] = "R";

constexpr int COMMON_PARTS = 3;
constexpr int MAX_ARGS = 3;
constexpr int MAX_PARTS = 6;

using $parts_array_t$ = $std$::$array$<$std$::$string_view$, MAX_PARTS>;

static $std$::$string$ #encode_write#(const $Event$& e)
    { return $std$::#format#(WRITE_FORMAT, e.args[0], e.args[1]); }

static $std$::$string$ #encode_pinmode#(const $Event$& e)
    { return $std$::#format#(PINMODE_FORMAT, e.args[0], e.args[1]); }

static $std$::$string$ #encode_reboot#(const $Event$& e)
    { return REBOOT_FORMAT; }

$std$::$string$ #encode_event#(const $Event$& event)
{
    $std$::$string$ commonPart = $std$::#format#(
        COMMON_PART_FORMAT, 
        event.timestampMicros, 
        event.localVirtualTime);

    switch (event.type) {
        case $Event$::$Type$::Write:
            return commonPart + #encode_write#(event);
        case $Event$::$Type$::PinMode:
            return commonPart + #encode_pinmode#(event);
        case $Event$::$Type$::Reboot:
            return commonPart + #encode_reboot#(event);
        default:
            return UNKNOWN_EVENT;
    }
}

bool #decode_int#($std$::$string_view$ str, int64_t& out) { ... }
bool #decode_int#($std$::$string_view$ str, int& out) { ... }
$parts_array_t$ #split#($std$::$string_view$ messageView) { ... }

bool #decode_event#(const $std$::$string$& message, $Event$& event)
{
    $parts_array_t$ parts = #split#(message);
    $std$::$array$<int, MAX_ARGS> args;

    for (int i = 0; i < MAX_ARGS; ++i) {
        if (!#decode_int#(parts[COMMON_PARTS + i], args[i])) break;
    }

    if (parts[2] == "W") {
        event = $Event$::#write#(args[0], args[1]);
    }
    else if (parts[2] == "P") {
        event = $Event$::#set_pinmode#(args[0], ($PinMode$)args[1]);
    }
    else {
        return false;
    }
    if (!#decode_int#(parts[0], event.timestampMicros) || 
        !#decode_int#(parts[1], event.localVirtualTime)) {
        return false;
    }
    return true;
}
\end{lstlisting}

\section{TCP Server module}

\subsection{\texttt{TcpServer} class}

\begin{lstlisting}[language=csharp, caption={\texttt{TcpServer} class}]
class $TcpServer$ {
   public:
    using $handler_t$ = $ConnectionHandler$;
    using $logger_ptr$ = $std$::$shared_ptr$<logger::$ILogger$>;

    $TcpServer$(
        uint16_t port, 
        $handler_t$::$message_funct$ messageHandler, 
        $logger_ptr$ logger);

    void #launch#();
    void #post_message#(const $std$::$string$& message);
    bool #is_connected#() const;

   private:
    $boost$::$asio$::$io_context$ _ioContext;
    $boost$::$asio$::$ip$::$tcp$::$endpoint$ _endpoint;
    $boost$::$asio$::$ip$::$tcp$::$acceptor$ _acceptor;
    uint16_t _port;

    $std$::$shared_ptr$<$handler_t$> _connectionHandler;
    $std$::$mutex$ _connectionMutex;
    $std$::$condition_variable$ _connectionCv;
    $handler_t$::$message_funct$ _messageHandler;
    $logger_ptr$ _logger;

    void #start_accepting#();
};
\end{lstlisting}

\subsection{\texttt{ConnectionHandler} class}

\begin{lstlisting}[language=csharp, caption={\texttt{ConnectionHandler} class}]
class $ConnectionHandler$ : 
    public $std$::$enable_shared_from_this$<$ConnectionHandler$> {
   public:
    static constexpr char MESSAGE_DELIMITER = ';';
    static constexpr size_t MAX_QUEUE_SIZE = 100;
    using $self_ptr$ = $std$::$shared_ptr$<$ConnectionHandler$>;
    using $logger_ptr$ = $std$::$shared_ptr$<logger::$ILogger$>;
    using $message_funct$ = $std$::$function$<void(const $std$::$string$&)>;

    $ConnectionHandler$(
        $boost$::$asio$::$io_context$& ioContext, 
        $message_funct$ messageHandler, 
        $logger_ptr$ logger);
    static $self_ptr$ #create#(
        $boost$::$asio$::$io_context$& ioContext, 
        $message_funct$ messageHandler, 
        $logger_ptr$ logger);

    void #handle#();
    void #post_message#(const $std$::$string$& message);
    $boost$::$asio$::$ip$::$tcp$::$socket$& #socket#() 
        { return _socket; }
    bool #is_connected#() const { return _isConnected; }

   private:
    $boost$::$asio$::$io_context$::$executor_type$ _executor;
    $boost$::$asio$::$ip$::$tcp$::$socket$ _socket;
    $std$::$string$ _incomingBuffer;
    $std$::$queue$<$std$::$string$> _outboxQueue;
    $std$::$mutex$ _queueMutex;
    $std$::$condition_variable$ _queueCv;
    size_t _connectionId;
    static size_t _nextConnectionId;
    bool _isConnected;
    $std$::$atomic$<bool> _isWriting;
    $message_funct$ _messageHandler;
    $logger_ptr$ _logger;

    void #handle_message#(
        $boost$::$system$::$error_code$ error, 
        size_t messageSize);
    void #write_from_queue#();
    void #disconnect#();
};
\end{lstlisting}

\section{Shared Memory module}

\subsection{\texttt{MemoryRegion} class}

\begin{lstlisting}[language=csharp, caption={\texttt{MemoryRegion} class}]
class $MemoryRegion$ {
   public:
    using $shmem_t$ = 
        $boost$::$interprocess$::$shared_memory_object$;
    using $shmem_region_t$ = 
        $boost$::$interprocess$::$mapped_region$;

    $MemoryRegion$(
        const $shmem_t$& shmemObject, 
        size_t offset, 
        size_t size);

    size_t #size#() { return _mappedRegion.#get_size#(); }
    uint8_t* #begin#() 
        { return (uint8_t*)_mappedRegion.#get_address#(); }
    uint8_t* #end#() { return #begin#() + #size#(); }

    template <typename T>
    T& #at#(size_t offset);
    uint8_t* #at#(size_t offset) { return #begin#() + offset; }
    size_t #offset#(void* ptr) 
        { return (uint8_t*)ptr - #begin#(); }

   private:
    $shmem_region_t$ _mappedRegion;
};
\end{lstlisting}

\subsection{\texttt{CircularBuffer} class}

\begin{lstlisting}[language=csharp, caption={\texttt{CircularBuffer} class}]
class $CircularBuffer$ {
   public:
    struct $iterator_t$;
    friend struct $iterator_t$;

    using $shmem_t$ = 
        $boost$::$interprocess$::$shared_memory_object$;

    static constexpr size_t PRODUCER_PTR_LOCATION = 0;
    static constexpr size_t CONSUMER_PTR_LOCATION = 8;
    static constexpr size_t BUFFER_LOCATION = 16;
    $CircularBuffer$() = default;
    $CircularBuffer$(
        const $shmem_t$& shmemObject, 
        size_t offset, 
        size_t size);

    bool #is_empty#();
    size_t #buffer_size#() 
        { return _memoryRegion->#size#() - BUFFER_LOCATION; }
    size_t #bytes_filled#();
    size_t #bytes_available#() 
        { return #buffer_size#() - #bytes_filled#(); }

    template <typename T>
        requires $std$::$is_trivially_copyable_v$<T>
              && (!$std$::$ranges$::$contiguous_range$<T>)
    bool #write#(const T& data);

    template <$std$::$ranges$::$contiguous_range$ TRange>
    bool #write#(const TRange& data);

    $std$::$vector$<uint8_t> #consume_until#(uint8_t delimeter);

   private:
    $std$::$unique_ptr$<$MemoryRegion$> _memoryRegion;

    $iterator_t$ #begin#();
    $iterator_t$ #end#();
    $iterator_t$ #at#(size_t offset);
    $iterator_t$ #producer_it#();
    $iterator_t$ #consumer_it#();
    void #write_producer#($iterator_t$ producer_it);
    void #write_consumer#($iterator_t$ consumer_it);
};
\end{lstlisting}

\subsection{\texttt{TwoWayBuffer} class}

\begin{lstlisting}[language=csharp, caption={\texttt{TwoWayBuffer} class}]
class $TwoWayBuffer$ {
   public:
    using $shmem_t$ = 
        $boost$::$interprocess$::$shared_memory_object$;

    $TwoWayBuffer$(const $std$::$string$& shmemName, size_t sizeInPages);
    ~$TwoWayBuffer$();

    template <typename T>
    bool #write#(const T& data);

    $std$::$vector$<uint8_t> #consume_until#(uint8_t delimeter);
    $CircularBuffer$& #producer#() { return _producer; }
    $CircularBuffer$& #consumer#() { return _consumer; }
    $std$::$string$& #name#() { return _name; }
    size_t #pages#() { return _pages; }
    size_t #size#();

   private:
    $std$::$string$ _name;
    size_t _pages;
    $shmem_t$ _shmemObject;
    $CircularBuffer$ _producer;
    $CircularBuffer$ _consumer;
};
\end{lstlisting}